\section{Fundamentos da Inteligência Artificial}

\textbf{OBJETIVO GERAL:}
Desenvolver competências relacionadas aos fundamentos científicos, matemáticos, computacionais e éticos da Inteligência Artificial, possibilitando ao estudante compreender os princípios que sustentam os sistemas inteligentes e sua aplicação em ambientes corporativos, industriais e tecnológicos.

\textbf{CARGA HORÁRIA:} 80 horas


\textbf{FUNÇÃO:} Compreender os fundamentos científicos e tecnológicos da Inteligência Artificial para aplicação em soluções computacionais.


\textbf{SUBFUNÇÕES:}
\begin{itemize}
\item Interpretar os fundamentos da Inteligência Artificial.
\item Relacionar conceitos matemáticos utilizados em modelos inteligentes.
\item Preparar ambientes computacionais para desenvolvimento de aplicações de IA.
\item Aplicar princípios éticos, legais e de governança relacionados ao desenvolvimento de Inteligência Artificial.
\end{itemize} 

\textbf{PADRÕES DE DESEMPENHO:}
\begin{itemize}
    \item Identificar problemas passíveis de solução por Inteligência Artificial.
    \item Diferenciar as principais áreas da IA.
    \item Reconhecer limitações dos modelos inteligentes.
    \item Interpretar conceitos fundamentais utilizados em Machine Learning.
    \item Aplicar princípios matemáticos básicos na interpretação de algoritmos.
    \item Configurar ambiente computacional para desenvolvimento de aplicações.
    \item Utilizar ferramentas de versionamento.
    \item Atuar observando princípios éticos e legais.
\end{itemize}

\textbf{CAPACIDADES TÉCNICAS:}

\textbf{Fundamentos}
\begin{itemize}
    \item Interpretar a evolução histórica da Inteligência Artificial.
    \item Diferenciar IA simbólica, estatística e generativa.
    \item Identificar aplicações da IA em diferentes setores econômicos.
    \item Reconhecer oportunidades de utilização da IA.
    \item Compreender limitações tecnológicas atuais.
\end{itemize}

\textbf{Matemática}
\begin{itemize}
    \item Interpretar vetores e matrizes.
    \item Aplicar conceitos de álgebra linear.
    \item Aplicar estatística descritiva.
    \item Interpretar distribuições estatísticas.
    \item Relacionar probabilidade e aprendizado de máquina.
\end{itemize}

\textbf{Computação}
\begin{itemize}
    \item Identificar arquiteturas computacionais.
    \item Compreender sistemas operacionais.
    \item Utilizar terminal Linux.
    \item Navegar em ambientes de desenvolvimento.
\end{itemize}

\textbf{Engenharia de Software}
\begin{itemize}
    \item Compreender versionamento.
    \item Utilizar repositórios Git.
    \item Interpretar documentação técnica.
    \item Utilizar ambientes colaborativos.
\end{itemize}

\textbf{CONHECIMENTOS:}

\textbf{História da IA}
\begin{itemize}
    \item Origem da Inteligência Artificial
    \item Conferência de Dartmouth
    \item Evolução histórica
    \item Inverno da IA
    \item Renascimento da IA
\end{itemize}

\textbf{Conceitos Fundamentais}
\begin{itemize}
    \item Inteligência Artificial
    \item Machine Learning
    \item Deep Learning
    \item IA Generativa
    \item IA Simbólica
    \item IA Estatística
\end{itemize}

\textbf{Áreas da IA}
\begin{itemize}
    \item NLP
    \item Visão Computacional
    \item Sistemas Especialistas
    \item Robótica
    \item Planejamento
    \item Otimização
    \item Agentes Inteligentes
\end{itemize}

\textbf{Álgebra Linear}
\begin{itemize}
    \item Escalares
    \item Vetores
    \item Matrizes
    \item Multiplicação Matricial
    \item Determinantes
\end{itemize}

\textbf{Estatística}
\begin{itemize}
    \item Média
    \item Mediana
    \item Moda
    \item Variância
    \item Desvio Padrão
\end{itemize}

\textbf{Probabilidade}
\begin{itemize}
    \item Eventos
    \item Distribuições
    \item Probabilidade Condicional
    \item Teorema de Bayes
\end{itemize}

\textbf{Sistemas Operacionais}
\begin{itemize}
    \item Linux
    \item Windows
    \item Terminal
    \item Shell
\end{itemize}

\textbf{Python}
\begin{itemize}
    \item Instalação
    \item Pip
    \item Virtualenv
    \item Conda
\end{itemize}

\textbf{Ética}
\begin{itemize}
    \item IA Responsável
    \item Transparência
    \item Viés Algorítmico
    \item Explainability
\end{itemize}

\textbf{Legislação}
\begin{itemize}
    \item LGPD
    \item Direitos Autorais
    \item Uso Ético de Dados
\end{itemize}

\textbf{Governança}
\begin{itemize}
    \item Auditoria
    \item Segurança
    \item Compliance
    \item Gestão de Riscos
\end{itemize}

\textbf{CAPACIDADES SOCIOEMOCIONAIS:}
\begin{itemize}
    \item Demonstrar disciplina na organização dos ambientes computacionais.
    \item Zelar pela documentação das atividades desenvolvidas.
    \item Demonstrar raciocínio lógico.
    \item Avaliar criticamente resultados obtidos.
    \item Formular hipóteses para resolução de problemas.
    \item Comunicar informações técnicas de forma clara.
    \item Produzir documentação objetiva.
    \item Compartilhar conhecimento com equipes.
    \item Colaborar em projetos coletivos.
    \item Compartilhar responsabilidades.
    \item Respeitar opiniões divergentes.
    \item Agir com responsabilidade no uso da Inteligência Artificial.
    \item Respeitar privacidade dos dados.
    \item Observar princípios de IA Responsável.
    \item Reconhecer riscos associados ao uso inadequado da IA.
\end{itemize}

\textbf{AMBIENTES PEDAGÓGICOS:}
\begin{itemize}
    \item Laboratório de Informática.
    \item Ambiente Virtual de Aprendizagem.
    \item Laboratório de Desenvolvimento de Software.
    \item Ambiente Cloud para execução de notebooks.
\end{itemize}

\textbf{EQUIPAMENTOS E RECURSOS:}

\textbf{Hardware}
\begin{itemize}
    \item Computadores com suporte à virtualização.
    \item Acesso à Internet.
    \item Projetor multimídia.
\end{itemize}

\textbf{Software}
\begin{itemize}
    \item Python
    \item VS Code
    \item Git
    \item GitHub Desktop
    \item Jupyter Notebook
    \item Google Colab
    \item Docker (introdução)
\end{itemize}
